\section{Actual whole division by the canonical reduced standard basis}

The earlier division constructs arbitrary monic vertex members of the actual
whole ideal. Those members need not have the same selected-centre orders as
the canonical reduced members. We therefore perform the actual recursion for
a specified monic family, then instantiate it with the internally constructed
canonical reduced basis. The construction follows the standard-basis division
in \cite{bierstone-milman03-desingularization-algorithms}, Section 6.3. This chapter does not infer centre compatibility
from an arbitrary monic choice or construct a semicoherent presentation.

\begin{theorem}[Whole division by an actual monic vertex family]
\label{hir:actual-monic-vertex-family-whole-division}
\source{bierstone-milman03-desingularization-algorithms}
\uses{hir:actual-formal-vertex-division-identity,hir:actual-formal-ideal-initial-diagram-upper}
\lean{Hironaka.HironakaActualReducedBasisDivision.ActualMonicVertexFamily,Hironaka.HironakaActualReducedBasisDivision.familyDivisionCoefficient,Hironaka.HironakaActualReducedBasisDivision.familyDivisionCoefficient_eq,Hironaka.HironakaActualReducedBasisDivision.familyFormalQuotient,Hironaka.HironakaActualReducedBasisDivision.familyFormalRemainder,Hironaka.HironakaActualReducedBasisDivision.familyFormalQuotient_support,Hironaka.HironakaActualReducedBasisDivision.familyFormalRemainder_coefficient_diagram,Hironaka.HironakaActualReducedBasisDivision.family_formal_vertex_division,Hironaka.HironakaActualReducedBasisDivision.familyFormalRemainder_eq_zero_of_mem,Hironaka.HironakaActualReducedBasisDivision.family_formal_ideal_member_division,Hironaka.HironakaActualReducedBasisDivision.family_formal_ideal_eq_span_vertex_series}
\leanok
Let $I\subset\mathbb C[[x]]$ have its genuine initial diagram and actual
vertices. For each actual vertex $a$, let $B_a\in I$ have coefficient one
at $a$ and zero coefficients strictly below $a$. Every formal $F$ has
internally constructed series $Q_a,R$ satisfying
\[
 F=\sum_a Q_aB_a+R,
\]
with the actual disjoint owner-part support for each $Q_a$ and zero
coefficients of $R$ on the entire actual diagram. If $F\in I$, then
$R=0$, and the $B_a$ generate the whole ideal.
\end{theorem}
\begin{proof}
Use the already constructed actual owner partition of the whole diagram.
At exponent $c$, define the correction coefficient by well-founded
recursion: subtract all finite antidiagonal products whose quotient
exponent plus its owner vertex is strictly below $c$. Put the correction
in that owner's quotient, or in the remainder if no owner exists.
For each antidiagonal product, the monic coefficient supplies the owner
term at $c$; the remaining terms are exactly the recursive tail. Terms
whose source vertex would lie above their second exponent vanish by the
assumed genuine lower coefficients. Summing all actual owners proves the
identity coefficient by coefficient. If $F\in I$, the finite product sum
lies in $I$, hence so does $R$. A nonzero such $R$ would have its genuine
initial exponent in the diagram, contradicting its vanishing there.
Whole ideal generation follows in both directions.
\end{proof}

\begin{theorem}[The same actual canonical normal representative]
\label{hir:actual-monic-family-canonical-remainder}
\uses{hir:actual-monic-vertex-family-whole-division,hir:actual-formal-normal-representative-unique}
\lean{Hironaka.HironakaActualReducedBasisDivision.familyFormalRemainder_eq_canonical}
\leanok
The actual remainder produced by any such family is exactly the previously
constructed canonical normal representative of the same class modulo $I$.
\end{theorem}
\begin{proof}
The whole division identity puts $F-R$ in $I$. Its actual remainder
coefficients vanish on the diagram. Apply uniqueness of the canonical
normal representative to these two already proved properties. No equality
of arbitrary chosen monic representatives is required.
\end{proof}

\begin{theorem}[Actual selected-centre filtration of the same quotients]
\label{hir:actual-monic-family-centre-power-division}
\uses{hir:actual-monic-vertex-family-whole-division,hir:actual-formal-coordinate-power-coefficient-kernel}
\lean{Hironaka.HironakaActualReducedBasisDivision.familyDivisionCoefficient_eq_zero_of_weight_lt,Hironaka.HironakaActualReducedBasisDivision.familyFormalQuotient_coefficient_weight_bound,Hironaka.HironakaActualReducedBasisDivision.familyFormalRemainder_coefficient_weight_bound,Hironaka.HironakaActualReducedBasisDivision.family_formal_ideal_member_weighted_division,Hironaka.HironakaActualReducedBasisDivision.familyFormalQuotient_mem_coordinate_power}
\leanok
Let $P_A$ be the actual coordinate ideal of a selected subset and write
$w_A(a)$ for its literal exponent weight. If $B_a\in P_A^{w_A(a)}$
for every actual vertex and $F\in I\cap P_A^q$, then the very same whole
division quotients satisfy $Q_a\in P_A^{q-w_A(a)}$.
\end{theorem}
\begin{proof}
The actual coordinate-power kernel translates the hypotheses into lower
coefficient vanishing. Induct along the very same well-founded recursion.
In a nonzero tail product, the second factor has weight at least the
owner vertex weight. Its first exponent plus the owner vertex therefore
has weight at most the total exponent. At a total weight below $q$ the
induction hypothesis kills its correction coefficient. Thus every
correction coefficient below $q$ vanishes. The quotient's actual owner
formula shifts this bound by the vertex weight. Translate it back through
the entire coordinate-power kernel. Power zero and empty coordinate blocks
are covered by the same argument.
\end{proof}

\begin{corollary}[Internally supplied canonical reduced family]
\label{hir:actual-reduced-basis-entire-filtered-division}
\uses{hir:actual-reduced-formal-vertex-members,hir:actual-monic-family-canonical-remainder,hir:actual-monic-family-centre-power-division}
\lean{Hironaka.HironakaActualReducedBasisDivision.actualReducedVertexFamily,Hironaka.HironakaActualReducedBasisDivision.actual_reduced_basis_entire_division,Hironaka.HironakaActualReducedBasisDivision.actual_reduced_basis_entire_filtered_division}
\leanok
Use the actual reduced members $G_a=X^a-R_I(X^a)$ of the same whole ideal.
They internally supply every monic-family hypothesis above. Hence there
is genuine whole division by these very reduced members, with the same
canonical remainder and the asserted centre-power quotient bounds whenever
their actual selected-centre powers hold.
\end{corollary}
\begin{proof}
The previous construction proves actual ideal membership, coefficient one
at the genuine vertex, and vanishing below it for every $G_a$. These are
exactly the three inputs to the actual recursion, with no supplied quotient,
normal form, generation certificate or filtered division. Substitute this
internally constructed family in the preceding theorems. Actual original
coefficient cosupports can supply the remaining centre-power hypotheses by
the separate coefficient-presentation producer; general presentation
construction and Bennett inequalities are still separate obligations.
\end{proof}
